\chapter{Design and Implementation}

\section{Overview of Artefact}
% Brief recap: IDR + dashboard
% Link back to DSR (Design -> instantiation)

The artefact developed in this study consists of a lightweight Integrated Data Repository (IDR) and associated data processing pipeline designed to support operational analysis of post-MDT oncology workflows. The implementation translates to the conceptual archetecture proposed in \ref{Chapter 3} into a functioning system capable of integrating data from multiple operational sources and producing am analytics-ready dataset. 

The system adopts a staged archetecture comprising: 
\begin{enumerate}
    \item Data ingestion from multiple NHS operational systems,
    \item A file-based staging layer,
    \item A transformation and integration layer, and
    \item A central analytical data store implemented in PostgreSQL. 
\end{enumerate}

This design reflects the ``general model'' of IDR archetectures, incorporating a staging layer to support extract-transform-load (ETL) processes and harmonisation accross heterogeneous data sources, as reccomended in healthcare informatics literature. 

\section{Data Sources and Ingestion}


\section{IDR Architecture}
% As built - not just proposed
%% Staging → warehouse → dashboard layer
%Actual technologies used (SQL Server? Power BI? Python ETL?)
%Data flow diagram (strongly recommended)
% Add explicit reflection:
%What changed from your proposed “target state IDR”?



\section{Data Engineering and ETL}
%Add:
%Source systems used (MDT, PAS, etc.)
%Data extraction approach
%Transformations:
%interval derivations (MDT→clinic etc.)
%breach-risk calculations
%Data quality handling (link to your section 8.2 in proposal)
%Include:
%Example transformation logic (pseudo-SQL is fine)
%Data dictionary excerpt

\section{Dashboard Design and Features}
%For each view:
%Post-MDT tracker
%Capacity panel
%Interval visualisations
%Explain:
%What it shows
%Why it exists (tie to DP1–DP6)
%How users interact with it

\section{Tools and Technologies}
%Example:
%SQL / ETL tooling
%Power BI / Tableau
%Data modelling approach
%Any scripting